\section*{ВВЕДЕНИЕ}
\addcontentsline{toc}{section}{ВВЕДЕНИЕ}

\textbf{Актуальность работы.} Корпоративные сети остаются устойчивой
целью несанкционированных воздействий, включающих эксфильтрацию
данных, латеральное перемещение, разведку, отказ в обслуживании и
эксплуатацию уязвимостей сетевых сервисов. Доля шифрованного трафика
в современных корпоративных инфраструктурах превышает 90\,\%, что
делает сигнатурный анализ полезной нагрузки малоэффективным.
Центр тяжести разработки систем обнаружения вторжений
(\eng{Intrusion Detection Systems, IDS}) переходит в сторону
аномального анализа метаданных потоков с применением методов
машинного и глубокого обучения. При этом существенной проблемой
остаётся воспроизводимая, формализованная оценка таких систем в
условиях, приближенных к реальной эксплуатации, в том числе при
дрейфе распределений и вариативности атакующих сценариев.

\textbf{Объект исследования} --- сетевой трафик корпоративной сети,
подверженной несанкционированным действиям и аномальной активности.

\textbf{Предмет исследования} --- методика обнаружения аномальной
активности и несанкционированных действий с использованием нейронных
сетей, а также методология её активного тестирования с помощью
программного агента ИИ-злоумышленника.

\textbf{Цель работы.} Разработка методики обнаружения аномальной
активности и несанкционированных действий в корпоративной сети на
основе нейронных сетей, включающей анализ сетевого трафика,
классификацию угроз и моделирование атак с помощью ИИ-злоумышленника,
генерирующего трафик в реальном времени для тестирования системы.

\textbf{Задачи исследования:}
\begin{enumerate}
    \item Изучение и систематизация существующих методов и технологий
    обнаружения аномалий и атак: статистические подходы, классические
    ML-методы, нейросетевые модели; обзор IDS/IPS- и NIDS-решений.
    \item Изучение методов моделирования атак с использованием ИИ и
    подходов к генерации синтетического трафика; определение роли
    ИИ-злоумышленника как агента переключения режимов
    \emph{норма/атака/дрейф}.
    \item Формулировка гипотезы исследования и формализация
    функциональных и нефункциональных требований к системе.
    \item Разработка и обоснование архитектуры нейронной сети,
    пригодной для последовательностей и окон сетевого трафика;
    определение входного представления и схемы алертирования.
    \item Сбор и подготовка данных сетевого трафика
    (\emph{UNSW-NB15, CICIDS2017}), включая препроцессинг,
    нормализацию, формирование скользящих окон и управление
    дисбалансом классов.
    \item Обучение нейросетевых моделей, настройка гиперпараметров,
    калибровка порогов под целевой эксплуатационный уровень ложных
    срабатываний, сравнение с baseline-моделями.
    \item Реализация ИИ-злоумышленника и комплексное тестирование
    прототипа в связке \emph{агент $\to$ детектор}.
    \item Реализация программного прототипа мониторинга и детекции
    в реальном времени с подсистемами интеграции с источниками
    сетевой телеметрии.
    \item Оценка эффективности, оптимизация и доработка прототипа,
    формирование выводов и направлений дальнейшего развития.
\end{enumerate}

\textbf{Гипотеза исследования.} Использование нейронных сетей для
анализа сетевого трафика позволяет с высокой точностью выявлять
несанкционированные действия и аномальную активность, проявляющиеся в
характерных сетевых паттернах (изменения интенсивности, нетипичные
направления, аномальные сессии и последовательности соединений).
Применение ИИ-злоумышленника в форме программного агента,
генерирующего трафик и атакующие сценарии в реальном времени (включая
сценарии с дрейфом распределений), обеспечивает полноту и
воспроизводимость тестирования, повышает устойчивость детектора и
снижает нагрузку на SOC за счёт уменьшения ложных срабатываний и
более точной приоритизации инцидентов.

\textbf{Методы исследования.} Систематический анализ литературы;
математическое моделирование (вероятностные и информационно-теоретические
меры дрейфа, формальные определения функций потерь и метрик);
обучение нейронных сетей и классических моделей в равных условиях;
статистический анализ результатов через bootstrap-доверительные
интервалы и парный критерий Уилкоксона.

\textbf{Научная новизна.} Новизна состоит в комбинации трёх
компонент: (а)~гибридной архитектуры CNN-LSTM для классификации
flow-окон, обоснованной теоретическим и экспериментальным
сравнением с десятью альтернативами в равных условиях;
(б)~ИИ-злоумышленника как воспроизводимого активного источника
тестовых сценариев на основе шаблонов, конечного автомата и
drift-инжектора, без использования генеративных моделей;
(в)~единого протокола экспериментов, включающего сравнение моделей в
равных условиях, анализ ошибок, калибровку порога под
эксплуатационный FPR, стресс-тестирование с агентом,
drift-устойчивость и эксплуатационную производительность.

\textbf{Практическая значимость.} Разработанный прототип реализован
на стеке Python\,/\,PyTorch с REST-интерфейсом на FastAPI и
визуальной панелью на Streamlit и может быть интегрирован в
подсистему сетевой аналитики корпоративного SOC. Зафиксированы
рекомендации по интеграции с источниками NetFlow\,/\,IPFIX,
калибровке порогов под допустимый FPR и оценке устойчивости к
дрейфу.

\textbf{Структура отчёта.} Основная часть отчёта состоит из четырёх
разделов, соответствующих четырём этапам, заявленным в постановке
задачи. \textbf{Этап~1} содержит теоретический анализ методов
обнаружения аномалий, обзор архитектур нейронных сетей, сравнение
датасетов NIDS, обзор методов моделирования атак с применением ИИ,
формализацию требований и гипотезы. \textbf{Этап~2} --- сравнительный
анализ архитектур нейронных сетей по опубликованным NIDS-результатам,
обоснование выбора целевой архитектуры, теоретическая модель данных
и pipeline, библиотека синтетических тестовых сценариев и протокол
экспериментов. \textbf{Этап~3} --- программная реализация прототипа
системы обнаружения аномалий, обучение и оптимизация моделей,
разработка ИИ-злоумышленника. \textbf{Этап~4} --- комплексное
тестирование прототипа, анализ эффективности, оптимизация и
доработка, оформление итоговых материалов; именно этому этапу уделено
наибольшее внимание (расчёт метрик, анализ ошибок, тесты с
ИИ-злоумышленником, оценка устойчивости к дрейфу и
производительности).
